\section{Finite windows of curves}
\label{sec:global-audit-no-go}

Numerical checks of BSD-type statements are made on finite lists of
curves.  This section records the elementary but easily forgotten fact
that a statement about all curves cannot be decided from a fixed finite
list.  The list used for examples in this paper is the five-curve window
\[
  W=\{11a1,\ 37a1,\ 121b1,\ 960d1,\ 571a1\}.
\]

\begin{theorem}[A finite window does not decide a universal statement]\label{thm:apparatus:finite-window-no-go}
Let $\mathcal C$ be the set of isomorphism classes of elliptic curves over
$\Q$, and consider assignments $\tau:\mathcal C\to\{\mathrm{true},
\mathrm{false}\}$, thought of as truth values of some property of curves.
Let $U(\tau)$ be the statement ``$\tau(E)=\mathrm{true}$ for every $E$''.
\begin{enumerate}
\item No function of the restriction $\tau|_W$ computes $U(\tau)$ for all
  assignments $\tau$: the assignment $\tau_1\equiv\mathrm{true}$ and the
  assignment $\tau_2$ that agrees with $\tau_1$ except at one curve
  $E_*\notin W$, where it is false, have the same restriction to $W$, but
  $U(\tau_1)$ is true and $U(\tau_2)$ is false.
\item In Schur form, let $\mathcal C(B)$ be the finite set of curves of
  conductor at most $B\ge960$, with $M(B)=\#\mathcal C(B)$ coordinates.
  The window source is the coordinate projection onto the five
  coordinates in $W$, and the residual against the identity readout on
  $\mathcal C(B)$ has trace $M(B)-5$, which tends to infinity with $B$.\footnote{Lean proves both parts for a model in which curves are indexed
by natural numbers and the window is the first five indices: the two
completions agree on the window, differ outside it, are separated by the
universal target, and are identified by every function of the window;
and the residual trace $M-5$ takes every natural-number value.  The model
says nothing about which truth assignments arise from actual properties
of curves.}
\end{enumerate}
\end{theorem}

\begin{proof}
(1) Any function of $\tau|_W$ takes the same value on $\tau_1$ and
$\tau_2$, while $U$ separates them.  (2) The trace of a coordinate
projection is the number of retained coordinates, so
$\tr(I_{M(B)}-P_W)=M(B)-5$.  There are infinitely many isomorphism classes
of elliptic curves over $\Q$ and finitely many of each conductor, so
$M(B)\to\infty$.
\end{proof}

The point of the statement is methodological.  A conclusion about all
curves needs an argument that reaches the curves outside the window, for
instance a theorem covering a family, or an exhaustion argument, or a
passage from density statements to pointwise ones.  The finite window by
itself supplies none of these.
